\chapter{函式}
\label{ch:function}

到目前為止，我們寫過的程式幾乎都只有一個 \texttt{main} 函式：所有變數、
所有邏輯都擠在同一層裡。 當程式變長、或某段邏輯需要重複使用時，把它
獨立寫成一個\textbf{函式}（function）會讓程式更容易閱讀、修改、測試。
一個函式就像一個小工具：給它固定形式的輸入（\textbf{參數}，parameter），
它做完一件明確的事，然後用 \texttt{return} 交回一個結果；中間怎麼實作，
呼叫它的人完全不需要知道。

這一章分四段推進：先學會寫有回傳值、無回傳值與多參數函式，並用純函式
對上 \texttt{void} 的設計差異；再從呼叫堆疊解釋整數與陣列參數為什麼行
為不同；接著談函式原型、標頭檔與多檔編譯；最後才看變參 \texttt{printf}
／\texttt{scanf}，以及遞迴的進階用法。 函式指標會在指標那一章正式展開。
堆疊與生命週期會在這一章建立直覺，到指標與遞迴章再串成完整圖像。

\section{定義有回傳值的函式}
\label{sec:func-return-value}

一個有回傳值的函式由幾個部分組成：\textbf{回傳型別}、\textbf{函式名}、
\textbf{參數列表}、\textbf{函式本體}。 本體內用 \texttt{return} 交回一個
值給呼叫者。

\begin{code}{c}{有回傳值的函式}
回傳型別 函式名(參數型別 參數名) {
    ...
    return 型別匹配的值;
}
\end{code}

範例 \thechapter.2 把「計算平方」獨立成 \texttt{square}；雖然
\texttt{n * n} 直接寫在 \texttt{main} 裡也可以，但包成函式之後，只要
看到 \texttt{square(x)} 這個名字，就不用重新想一遍它在做什麼。 呼叫
\texttt{square(n)} 時，\texttt{n} 的值被複製一份給 \texttt{square} 的
參數，計算完畢後回傳值再交給 \texttt{printf}。

\begin{example}{c}{(square.c) 計算平方的函式}
#include <stdio.h>

int square(int n) {
    return n * n;
}

int main(void) {
    int n;
    scanf("%d", &n);
    printf("%d\n", square(n));
}
\end{example}

\iosample{
7 \\
}{
49 \\
}

\subsection{main 函式}

\texttt{main} 也是一個函式，它的特殊之處是作為程式的\textbf{入口點}
（entry point）——作業系統啟動程式時，會呼叫 \texttt{main}。 \texttt{main}
回傳的整數值會交給作業系統，告知程式是否成功執行。 慣例上 \texttt{0}
表示成功，非 \texttt{0} 表示失敗。

\begin{code}{c}{main 的簽名}
int main(void) {
    // 程式邏輯
    return 0;  // 告知作業系統程式成功
}

// 或者接受命令行引數（進階用法，這裡先不詳述）
int main(int argc, char *argv[]) {
    // ...
}
\end{code}

\subsection{判斷閏年}

實際應用中，我們常把「某個明確條件判斷」包成一個函式。 例如判斷是否
為閏年：年份能被 4 整除但不能被 100 整除，或者能被 400 整除。

\begin{example}{c}{(leap-year.c) 判斷閏年的函式}
#include <stdio.h>

// 回傳 1 表示閏年，0 表示平年
int isLeapYear(int year) {
    if (year % 400 == 0) {
        return 1;  // 能被 400 整除，是閏年
    }
    if (year % 100 == 0) {
        return 0;  // 能被 100 整除但不能被 400 整除，是平年
    }
    if (year % 4 == 0) {
        return 1;  // 能被 4 整除但不能被 100 整除，是閏年
    }
    return 0;  // 其他都是平年
}

int main(void) {
    int year;
    scanf("%d", &year);
    printf("%d\n", isLeapYear(year));
}
\end{example}

\iosample{
2024 \\
}{
1 \\
}

輸入 \texttt{2024}，按第三個條件 \texttt{2024 \% 4 == 0}（且不被 100
整除），回傳 \texttt{1}。 輸入 \texttt{2000} 則按第一個條件回傳 \texttt{1}。

\keypoint{\texttt{return} 一執行就立刻結束這次函式呼叫，後面的程式碼
不會再跑。 用連續的 \texttt{if} 搭配提早 \texttt{return}，通常比複雜的
\texttt{if-else} 鏈更平坦好讀。}

\begin{drills}[label=drills:func-return-value]
\item \texttt{square} 函式呼叫 \texttt{square(n)} 時，\texttt{n} 的值
是直接傳進去，還是複製一份給參數？
\item \texttt{isLeapYear(2000)} 會先命中哪一個 \texttt{if}？回傳值是
多少？
\item 把 \texttt{isLeapYear} 裡三個 \texttt{if} 的判斷順序打亂（先判斷
\% 4，再判斷 \% 400），對 2000 這個輸入會不會算錯？為什麼？
\end{drills}

\section{定義無回傳值的函式}
\label{sec:func-void}

\texttt{void} 型別表示函式不回傳任何值。 這類函式通常用來執行某個動作
或修改某些狀態——比如印出結果、修改陣列、更新全域變數。

\begin{code}{c}{無回傳值的函式}
void 函式名(參數型別 參數名) {
    // 執行動作
    // 沒有 return 值（或 return; 後面不跟值）
}
\end{code}

範例~\ref{ex:print-digits} 寫一個 \texttt{void} 函式，把一個整數的各位數字
分別印在不同行。 這個函式不回傳任何計算結果，只是有副作用（把值印出來）。

\begin{example}[label=ex:print-digits]{c}{(print-digits.c) 將各位數印在不同行}
#include <stdio.h>

// 無回傳值的函式：只負責輸出，不計算結果
void printDigits(int n) {
    if (n < 0) n = -n;  // 處理負數

    if (n == 0) {
        printf("0\n");
        return;
    }

    // 遞迴式地從最高位開始印
    if (n >= 10) {
        printDigits(n / 10);
    }
    printf("%d\n", n % 10);
}

int main(void) {
    int num;
    scanf("%d", &num);
    printDigits(num);
}
\end{example}

\iosample{
12345 \\
}{
1 \\
2 \\
3 \\
4 \\
5 \\
}

\texttt{printDigits(12345)} 遞迴呼叫自己，依序把 \texttt{12345/10 = 1234}
傳進去，最後回到 \texttt{1}；再回溯時逐一印出 \texttt{2, 3, 4, 5}。

另一個 \texttt{void} 函式的例子是初始化或填充陣列：

\begin{code}{c}{用 void 函式初始化陣列}
void fillZeros(int arr[], int n) {
    for (int i = 0; i < n; i++) {
        arr[i] = 0;  // 副作用：修改陣列
    }
}
\end{code}

這類函式的副作用（修改陣列）是預期的，所以函式名往往暗示「會改東西」
（如 \texttt{fillZeros}、\texttt{reverse}、\texttt{sort}）。

\keypoint{\texttt{void} 函式沒有回傳值，呼叫它是為了它的副作用（印出
結果、修改陣列、改變全域狀態等）。 在函式名上要講清楚是否會有副作用。}

\begin{drills}[label=drills:func-void]
\item \texttt{printDigits} 沒有 \texttt{return} 值，那它是靠什麼讓
數字出現在螢幕上？
\item \texttt{fillZeros(arr, n)} 呼叫結束後，\texttt{main} 裡原本的
陣列內容會不會被改變？
\item 一個函式如果同時「有副作用」又「有回傳值」，合理嗎？舉一個你覺
得可以接受的例子。
\end{drills}

\section{多個參數的函式}
\label{sec:func-params}

函式可以接多個參數。 呼叫時，每個引數各複製一份值給對應的參數，彼此
獨立。

\begin{code}{c}{多參數函式}
回傳型別 函式名(型別1 參數1, 型別2 參數2, ...) {
    return 根據參數計算的結果;
}
\end{code}

\subsection{計算天數}

一個常見的例子是把「年、月、日」轉換成該年的第幾天。

\begin{example}{c}{(day-of-year.c) 計算日期是一年中的第幾天}
#include <stdio.h>

// 給定年月日，回傳該年第幾天
int dayOfYear(int year, int month, int day) {
    // 各月的天數
    int daysInMonth[] = {31, 28, 31, 30, 31, 30, 31, 31, 30, 31, 30, 31};

    // 處理閏年
    if (year % 400 == 0 || (year % 4 == 0 && year % 100 != 0)) {
        daysInMonth[1] = 29;
    }

    // 加總前面各月的天數
    int result = 0;
    for (int m = 0; m < month - 1; m++) {
        result += daysInMonth[m];
    }

    // 加上這個月的日期
    result += day;
    return result;
}

int main(void) {
    int y, m, d;
    scanf("%d %d %d", &y, &m, &d);
    printf("%d\n", dayOfYear(y, m, d));
}
\end{example}

\iosample{
2024 3 1 \\
}{
61 \\
}

\subsection{印出連續整數序列}

另一個例子是印出從 \texttt{start} 到 \texttt{end} 的連續整數。

\begin{code}{c}{印出一個數字範圍}
void printRange(int start, int end) {
    for (int i = start; i <= end; i++) {
        if (i > start) printf(" ");
        printf("%d", i);
    }
    printf("\n");
}
\end{code}

呼叫 \texttt{printRange(3, 7)} 會印出 \texttt{3 4 5 6 7}。

\subsection{算式判斷}

一些題目要求判斷一個算式是否成立。 例如判斷三個邊長是否能組成直角三角形。

\begin{example}{c}{(right-triangle.c) 判斷是否為直角三角形}
#include <stdio.h>

// 判斷三邊能否組成直角三角形
int isRightTriangle(int a, int b, int c) {
    // 找最長邊
    int max = a;
    if (b > max) max = b;
    if (c > max) max = c;

    // 檢查勾股定理：最長邊的平方 = 其他兩邊平方和
    if (a == max) {
        return a * a == b * b + c * c;
    }
    if (b == max) {
        return b * b == a * a + c * c;
    }
    return c * c == a * a + b * b;
}

int main(void) {
    int a, b, c;
    scanf("%d %d %d", &a, &b, &c);
    printf("%d\n", isRightTriangle(a, b, c));
}
\end{example}

\iosample{
3 4 5 \\
}{
1 \\
}

\subsection{比較判斷}

判斷兩個值的大小或相等關係也很常見。

\begin{code}{c}{比較函式}
int compare(int a, int b) {
    if (a < b) return -1;
    if (a > b) return 1;
    return 0;
}

// 或者直接回傳布林結果
int isGreater(int a, int b) {
    return a > b;
}
\end{code}

\keypoint{多參數函式的每個參數在呼叫時各拿到一份獨立的複本；參
數間互不干擾。}

\begin{drills}[label=drills:func-params]
\item \texttt{dayOfYear(2024, 3, 1)} 為什麼要先處理閏年，才能算出正確
的答案？
\item \texttt{printRange(3, 7)} 印出的最後一個數字是幾？
\item \texttt{isRightTriangle} 為什麼要先找出最長邊，才能檢查勾股定理？
\end{drills}

\section{函式的優點}
\label{sec:func-benefits}

到目前為止，我們已經寫過有回傳值的 \texttt{square}、\texttt{dayOfYear}，
也有只負責輸出的 \texttt{printDigits}。 這一節整理函式為什麼值得寫，以及
\texttt{void} 與有回傳值函式在設計上該怎麼分工。

\subsection{純函式}

一個\textbf{純函式}（pure function）給定相同的輸入，\textbf{永遠會回傳相同
的結果}，而且不會改變程式中其他地方的狀態（如全域變數、輸入輸出等）。
例如 \texttt{square(5)} 在哪裡呼叫、呼叫幾次，都一定回傳 \texttt{25}。

相反地，一個\textbf{有副作用的函式}（function with side effects）不只回傳
結果，還會「順便」改變其他東西——修改陣列、改寫全域變數、印出畫面等。 例如
\texttt{printDigits} 只用於印出；本節後面的 \texttt{incrementAndReturn} 還會改全域變數。

\begin{example}{c}{(pure-vs-impure.c) 純函式 vs 有副作用的函式}
#include <stdio.h>

// 純函式：給定 n，永遠回傳同一個值
int square(int n) {
    return n * n;
}

// 純函式：判斷是否偶數，不改變任何狀態
int isEven(int n) {
    return n % 2 == 0;
}

// 有副作用：修改全域變數
int count = 0;
int incrementAndReturn(int n) {
    count++;  // 副作用：修改全域變數
    return n + 1;
}

int main(void) {
    printf("%d\n", square(5));    // 純函式，多次結果同
    printf("%d\n", square(5));    // 輸出：25, 25

    printf("%d %d\n", incrementAndReturn(5), count);  // 輸出：6 1
    printf("%d %d\n", incrementAndReturn(5), count);  // 輸出：6 2
}
\end{example}

\iosample{
}{
25 \\
25 \\
6 1 \\
6 2 \\
}

同樣呼叫 \texttt{incrementAndReturn(5)} 兩次，第二次執行時全域變數
\texttt{count} 已經從 \texttt{1} 變成 \texttt{2}，結果不再「純粹」。

\subsection{純函式的優點}

\textbf{易於測試}：純函式不依賴外部狀態，測試時只需檢查輸入和輸出是否
對應。 有副作用的函式往往需要檢查「執行前後，哪些全域變數改了」，測試
更複雜。

\textbf{易於組合}：寫 \texttt{minOfThree(a, b, c)} 時，可以安心呼叫
\texttt{minOfTwo} 兩次，因為知道它不會改變 \texttt{a}、\texttt{b}、
\texttt{c} 的值。

\textbf{易於並行執行}：純函式的呼叫之間沒有隱藏的依賴，多個呼叫可以
同時進行而不會衝突。

\begin{example}{c}{(compose-functions.c) 組合純函式}
#include <stdio.h>

int min(int a, int b) {
    return a < b ? a : b;
}

int minOfThree(int a, int b, int c) {
    return min(min(a, b), c);
}

int main(void) {
    printf("%d\n", minOfThree(5, 2, 9));  // 輸出 2
}
\end{example}

\iosample{
}{
2 \\
}

\subsection{模組化與重用}

把邏輯獨立成函式的另一個好處是\textbf{重用}。 一旦寫好 \texttt{gcd}（最大
公因數），就能在求最小公倍數、約分、題目檢查等多個地方使用，不用重複寫。

\begin{example}{c}{(lcm.c) 用輔助函式組合邏輯}
#include <stdio.h>

// 歐幾里得演算法
int gcd(int a, int b) {
    while (b != 0) {
        int t = b;
        b = a % b;
        a = t;
    }
    return a;
}

// 用 gcd 組合出最小公倍數
int lcm(int a, int b) {
    return a / gcd(a, b) * b;
}

int main(void) {
    int n;
    scanf("%d", &n);

    int result;
    scanf("%d", &result);
    for (int i = 1; i < n; i++) {
        int x;
        scanf("%d", &x);
        result = lcm(result, x);
    }

    printf("%d\n", result);
}
\end{example}

\iosample{
3 \\
4 6 8 \\
}{
24 \\
}

\texttt{main} 自己完全沒有寫任何輾轉相除的邏輯，只是把資料交給
\texttt{lcm}；\texttt{lcm} 也不需要知道歐幾里得演算法的細節，只要相信
\texttt{gcd} 正確即可。

\subsection{控制複雜度}

當函式複雜時，可以再把它拆成更小的輔助函式。 這樣每層都只負責一個明確
的任務，整個程式讀起來就像在讀高階的敘述，而非深陷細節。

\subsection{修改與除錯}

有 bug 時，如果邏輯集中在一個函式裡，修改一個地方就能解決所有問題；如果
邏輯散落在 \texttt{main} 的各個角落，bug 修改起來就容易遺漏。

\keypoint{純函式好測試、好組合；有副作用的函式要在名字和註解裡講清楚。
優先寫純函式，只在必要時才接受副作用。}

\begin{drills}[label=drills:func-benefits]
\item \texttt{incrementAndReturn} 是純函式嗎？為什麼？
\item 為什麼「純函式易於組合」？用 \texttt{minOfThree} 呼叫
\texttt{min} 兩次來說明。
\item \texttt{lcm} 函式裡完全沒寫輾轉相除的邏輯，這是靠什麼原則做到的？
\end{drills}

\section{呼叫堆疊與區域變數的生命週期}
\label{sec:func-callstack}

上一節區分純函式與有副作用的函式時，我們強調好的設計盡量不改呼叫者
的狀態。 但就算函式看起來想改參數，在 C 裡通常也改不到 \texttt{main}
裡原本的整數變數。 原因在\textbf{呼叫堆疊}（call stack）：函式裡宣告
的變數（包含參數）都是\textbf{區域變數}（local variable），只存在於這
一次函式呼叫所佔用的 stack frame 裡；函式 \texttt{return} 後，這些變
數占用的記憶體就釋放，名字也隨之失效。 \texttt{main} 裡的 \texttt{int x}
跟某個函式裡的 \texttt{int x} 即使名字相同，也是兩塊完全不相干的記憶體。

\begin{center}
\begin{tikzpicture}
    \node[cell, minimum width=2.8cm, minimum height=0.75cm] (squareTop) {返回位址};
    \node[hlcell, minimum width=2.8cm, minimum height=0.75cm, below=0pt of squareTop] (squareLocals) {區域變數 \texttt{n}};
    \node[cell, minimum width=2.8cm, minimum height=0.75cm, below=0.75cm of squareLocals] (mainTop) {返回位址};
    \node[hlcell, minimum width=2.8cm, minimum height=0.75cm, below=0pt of mainTop] (mainLocals) {區域變數 \texttt{x}, \texttt{y}};

    \node[idxlbl, anchor=east, xshift=-0.15cm] at ($(squareTop.west)!0.5!(squareLocals.west)$) {\texttt{square}};
    \node[idxlbl, anchor=east, xshift=-0.15cm] at ($(mainTop.west)!0.5!(mainLocals.west)$) {\texttt{main}};

    \coordinate (squareArrowStart) at ($(squareLocals.east)+(0.4,0)$);
    \coordinate (squareArrowEnd) at ($(squareLocals.east)+(1.2,0)$);
    \draw[diagarrow] (squareArrowStart) -- (squareArrowEnd);
    \node[idxlbl, anchor=west, align=left, right=0.05cm of squareArrowEnd] {目前執行\\這一層};

    \node[idxlbl, anchor=west, align=left] at ($(mainTop.east)+(1.25,0)$) {暫停，等上面\\做完再回來};
\end{tikzpicture}
\diagcaption{呼叫 \texttt{square(n)} 時，堆疊頂端多一層 \texttt{square}
的 frame（淺灰是這次呼叫專屬的區域變數）；\texttt{square} 返回後，這一
層被整個彈出，\texttt{main} 的 \texttt{x}、\texttt{y} 仍在下面那一層
等待繼續執行。}
\label{fig:ch08-call-stack-square}
\end{center}

片語 \thechapter.16 用佔位符概括「一次函式呼叫」在堆疊上發生的事：編譯
器為這次呼叫準備一層 frame，把引數的值複製進參數，跳去執行函式本體，最
後把回傳值交給呼叫者，並銷毀這一層 frame。

\begin{code}{c}{一次函式呼叫在堆疊上的生命週期（概念）}
// 呼叫者這一層的 frame 仍在堆疊上，只是暫停
結果 = 被呼叫函式(引數的值);   // 引數「值」被複製一份給參數

// 被呼叫函式本體執行期間：
//   1. 在堆疊頂端建立新的 frame
//   2. 參數與函式內宣告的區域變數都活在這層 frame 裡
//   3. return 時銷毀這層 frame，把回傳值留給呼叫者
\end{code}

這個「參數是複製一份值進去」的規則，在 C 裡叫做\textbf{傳值呼叫}
（call by value）。 範例 \thechapter.17 寫了一個看起來像要交換兩個整數
的 \texttt{swap} 函式，但執行後 \texttt{main} 裡的 \texttt{x}、
\texttt{y} 完全沒變。 因為 \texttt{swap} 交換的是它自己 frame 裡那兩
份複本，跟 \texttt{main} 的 \texttt{x}、\texttt{y} 無關。

\begin{example}{c}{(swap-params.c) 在函式裡改參數，改不到呼叫者的變數}
#include <stdio.h>

void swap(int a, int b) {
    int t = a;
    a = b;
    b = t;
}

int main(void) {
    int x = 3, y = 7;
    swap(x, y);
    printf("%d %d\n", x, y);
}
\end{example}

\iosample{}{
3 7 \\
}

用 \texttt{x = 3}、\texttt{y = 7} 走一遍：\texttt{swap(x, y)} 把
\texttt{3} 複製給參數 \texttt{a}、\texttt{7} 複製給 \texttt{b}，在
\texttt{swap} 的 frame 裡 \texttt{a} 與 \texttt{b} 確實互換了，但
\texttt{main} 的 \texttt{x}、\texttt{y} 從頭到尾都在另一層 frame，值
分別仍是 \texttt{3} 和 \texttt{7}。 如果題目真的要在函式裡改到呼叫者
的變數，需要指標（下一章才會正式教）；現階段只要記住：一般整數參
數傳進去的是「值的複本」。

\begin{center}
\begin{tikzpicture}
    \node[hlcell2, minimum width=1.1cm] (mx) {3};
    \node[hlcell2, minimum width=1.1cm, right=0.2cm of mx] (my) {7};
    \node[idxlbl, left=0.3cm of mx] {\texttt{main}};
    \node[hlcell, minimum width=1.1cm, right=1.1cm of my] (pa) {3};
    \node[hlcell, minimum width=1.1cm, right=0.2cm of pa] (pb) {7};
    \node[idxlbl, left=0.3cm of pa] {\texttt{swap}};
    \draw[diagarrow, black!45] (mx.east) -- (pa.west);
    \draw[diagarrow, black!45] (my.east) -- (pb.west);
    \node[idxlbl, align=left, right=0.6cm of pb] {呼叫當下\\各複製一份};
\end{tikzpicture}
\diagcaption{\texttt{main} 的 \texttt{x}、\texttt{y}（深灰）與
\texttt{swap} 的參數 \texttt{a}、\texttt{b}（淺灰）是四個獨立位置；
在 \texttt{swap} 裡改 \texttt{a}、\texttt{b} 不會回寫到 \texttt{x}、
\texttt{y}。}
\label{fig:ch08-stack-frame-scope}
\end{center}

\keypoint{區域變數（含參數）只活在這次函式呼叫的 stack frame 裡；
函式 \texttt{return} 後，這些變數占用的記憶體就釋放，名字也隨之失效。
不同函式裡同名變數互不影響。}

\begin{drills}[label=drills:func-callstack]
\item \texttt{swap(x, y)} 執行完後，\texttt{main} 裡的 \texttt{x}、
\texttt{y} 為什麼沒有互換？
\item 函式裡宣告的區域變數，什麼時候會被釋放？
\item \texttt{main} 裡有一個 \texttt{int x}，某個函式裡也宣告了
\texttt{int x}，這是同一塊記憶體嗎？
\end{drills}

\section{參數傳遞}
\label{sec:func-passing}

上一節的 \texttt{swap} 已經示範：整數參數傳進去的是複本。 這一節把
整數、陣列、二維陣列分開講清楚：整數型別傳的是值的複本；陣列和結構體
有特殊的處理方式。

\subsection{整數參數傳遞}

整數參數遵循\textbf{傳值呼叫}（call by value）：引數的值被複製一份進
stack frame，函式內對參數的修改改不到呼叫者的變數。

\begin{code}{c}{整數傳值示例}
void increment(int n) {
    n++;  // 只改這個 frame 裡的副本
}

int main(void) {
    int x = 5;
    increment(x);
    printf("%d\n", x);  // 仍是 5，不是 6
}
\end{code}

這是一種\textbf{保護}：寫函式的人不用擔心參數會洩漏回去。 但如果題目要
求在函式內真的改到呼叫者的整數，就得用指標（下一章的內容）。

\subsection{一維陣列參數傳遞}

陣列是例外：傳給函式的\textbf{不是整份陣列的複本}，而是陣列的起始位址
（即第一個元素的指標）。 函式對陣列元素做的修改會直接反映到呼叫者的
原陣列上。

\begin{example}{c}{(modify-array.c) 陣列參數會被修改}
#include <stdio.h>

// 計算陣列裡偶數的個數，並把偶數改成 0
int countEvensAndNullify(int arr[], int n) {
    int count = 0;
    for (int i = 0; i < n; i++) {
        if (arr[i] % 2 == 0) {
            count++;
            arr[i] = 0;  // 副作用：直接改陣列
        }
    }
    return count;
}

int main(void) {
    int nums[] = {1, 2, 3, 4, 5, 6};
    int n = 6;

    int evenCount = countEvensAndNullify(nums, n);

    printf("偶數個數：%d\n", evenCount);
    printf("修改後的陣列：");
    for (int i = 0; i < n; i++) {
        if (i) printf(" ");
        printf("%d", nums[i]);
    }
    printf("\n");
}
\end{example}

\iosample{
}{
偶數個數：3 \\
修改後的陣列：1 0 3 0 5 0 \\
}

陣列參數的注意事項：
\begin{itemize}
\item 函式簽名中 \texttt{int arr[]} 看起來像陣列，但實際上只是一個指標。
\item 函式無法直接知道陣列大小，所以通常需要另外傳一個 \texttt{int n}。
\item 函式對陣列的修改是真實的、會反映到呼叫者的原陣列。
\end{itemize}

\subsection{多維陣列參數傳遞}

多維陣列傳進函式時，第一維可以省略，但之後的維度必須明確指定。

\begin{example}{c}{(matrix-operation.c) 二維陣列參數}
#include <stdio.h>

// 把 src 順時針旋轉 90 度，寫進另一塊記憶體 dst
void rotate90(int n, int src[n][n], int dst[n][n]) {
    for (int i = 0; i < n; i++) {
        for (int j = 0; j < n; j++) {
            dst[j][n - 1 - i] = src[i][j];
        }
    }
}

int main(void) {
    int n;
    scanf("%d", &n);

    int a[n][n], b[n][n];
    for (int i = 0; i < n; i++) {
        for (int j = 0; j < n; j++) {
            scanf("%d", &a[i][j]);
        }
    }

    rotate90(n, a, b);

    for (int i = 0; i < n; i++) {
        for (int j = 0; j < n; j++) {
            if (j) printf(" ");
            printf("%d", b[i][j]);
        }
        printf("\n");
    }
}
\end{example}

\iosample{
2 \\
1 2 \\
3 4 \\
}{
3 1 \\
4 2 \\
}

在 \texttt{rotate90(n, a, b)} 中，\texttt{a} 和 \texttt{b} 傳的是位址，
函式裡對 \texttt{dst[i][j]} 的寫入會直接改到 \texttt{main} 的 \texttt{b}
陣列。 這是多維陣列參數和一維陣列參數一樣的行為。

\keypoint{陣列參數傳的是位址而非複本；函式對陣列的修改會直接影響原陣列。
傳入多維陣列時，第一維可省略，但後續維度必須明確。}

\begin{drills}[label=drills:func-passing]
\item 整數參數與陣列參數傳給函式的分別是什麼？
\item \texttt{countEvensAndNullify(nums, n)} 執行完後，\texttt{main}
裡的 \texttt{nums} 陣列內容會不會改變？
\item 為什麼陣列參數通常還要另外傳一個 \texttt{int n}？
\end{drills}

\section{常見陷阱與最佳實踐}
\label{sec:func-pitfalls}

前兩節講過整數參數改不到原變數、陣列參數卻能改到元素、以及函式本身不知道
陣列有多大。 這一節把新手最常踩的三個坑集中列出，方便對照。

\subsection{陷阱 1：整數參數傳值，改不到原變數}

\begin{code}{c}{傳值的限制}
void tryToIncrement(int n) {
    n++;  // 只改副本
}

int main(void) {
    int x = 5;
    tryToIncrement(x);
    printf("%d\n", x);  // 仍是 5
}
\end{code}

這是設計特性，不是 bug。 要改呼叫者的變數，得用指標。

\subsection{陷阱 2：陣列參數與整數參數行為不同}

\begin{code}{c}{對比：整數 vs 陣列}
void modifyInt(int n) {
    n = 999;  // 改不到
}

void modifyArray(int arr[], int n) {
    arr[0] = 999;  // 能改到
}
\end{code}

新手往往忽略這個區別。 同一個道理——陣列傳的是位址——但機制與整數完全不同。

\subsection{陷阱 3：忘記陣列參數需要傳大小}

\begin{code}{c}{陷阱：函式不知道陣列有多大}
void sum(int arr[]) {
    // 無法知道 arr 有多少個元素！
    // int size = sizeof(arr);  // 這會給出指標的大小，不是陣列大小
}
\end{code}

總是一起傳陣列和它的大小：

\begin{code}{c}{正確做法}
int sum(int arr[], int n) {
    int total = 0;
    for (int i = 0; i < n; i++) {
        total += arr[i];
    }
    return total;
}
\end{code}

\subsection{最佳實踐}

\begin{itemize}
\item \textbf{優先寫純函式}：如果不需要副作用，就回傳結果；不要修改參數或
改全域變數。

\item \textbf{講清楚副作用}：若函式必須修改陣列，函式名要暗示這一點
（\texttt{reverse}、\texttt{fillWithZeros}），在註解裡也要說明。

\item \textbf{檢查回傳值}：\texttt{scanf}、\texttt{fopen} 等函式回傳狀態碼，
呼叫時要檢查。

\item \textbf{用函式原型}：大檔案裡，在前面集中寫原型；小檔案可以把定義
寫在前面。

\item \textbf{避免全域變數}：如果可以用參數和回傳值，就不用全域變數；
全域狀態容易導致副作用難預測。

\item \textbf{陣列總是要傳大小}：沒有例外。
\end{itemize}

\begin{drills}[label=drills:func-pitfalls]
\item 為什麼 \texttt{void modifyInt(int n) \{ n = 999; \}} 改不到呼叫者
的變數？
\item \texttt{sizeof(arr)} 在函式參數 \texttt{int arr[]} 裡量到的是什
麼？為什麼不能拿它當陣列長度？
\item 陣列參數需要修改時，函式名該不該暗示這件事？舉一個好名字和一個
壞名字。
\end{drills}

\section{呼叫函式}
\label{sec:func-declare}

到目前為止，函式定義都寫在同一個 \texttt{.c} 檔裡，而且多半放在
\texttt{main} 前面或後面。 程式變大、拆成多個檔案時，編譯器仍然必須在
每一次呼叫前知道函式「長什麼樣子」：它叫什麼名字、接什麼參數、回傳什麼
型別。 這就是 C 語言\textbf{宣告}與\textbf{定義}分離的由來。 有三種常見
方式告訴編譯器這些資訊。

\subsection{定義在前}

最簡單的方式是把函式定義寫在 \texttt{main} 之前。 這樣編譯器掃描到
\texttt{main} 時，已經看過函式的完整定義，知道 \texttt{square(n)} 存在
且回傳 \texttt{int}。

\begin{code}{c}{函式定義在 main 前面}
#include <stdio.h>

// 定義在前
int square(int n) {
    return n * n;
}

// 呼叫在後
int main(void) {
    int x;
    scanf("%d", &x);
    printf("%d\n", square(x));
}
\end{code}

\subsection{函式原型}

當檔案很大、或想把主要邏輯寫在 \texttt{main} 時，可以在前面先寫一個
\textbf{函式原型}（function prototype）——只告訴編譯器函式的簽名，不寫本體。

\begin{code}{c}{用函式原型告訴編譯器函式簽名}
#include <stdio.h>

// 函式原型（只有簽名，沒有本體）
int square(int n);

int main(void) {
    int x;
    scanf("%d", &x);
    printf("%d\n", square(x));
}

// 函式定義放在後面
int square(int n) {
    return n * n;
}
\end{code}

函式原型必須放在函式被呼叫之前，通常在檔案頂部集中寫。 編譯器見到
\texttt{square(x)} 時，會查詢前面宣告的原型，確認簽名正確。

\subsection{標頭檔}

當程式擴展到多個檔案時，用\textbf{標頭檔}（header file）是標準做法。
標頭檔（\texttt{.h}）裡放的是函式原型和常用的型別定義，源檔（\texttt{.c}）
裡放的是實作。 例如標準庫的 \texttt{stdio.h} 就是一個標頭檔，裡面宣告了
\texttt{printf}、\texttt{scanf} 等函式的原型；實際的實作在系統的標準庫
文件裡。

\begin{code}{c}{標頭檔的結構}
// math.h - 公開介面
int square(int n);
int add(int a, int b);

// math.c - 實作
#include "math.h"
int square(int n) {
    return n * n;
}
int add(int a, int b) {
    return a + b;
}

// main.c - 呼叫
#include "math.h"
int main(void) {
    printf("%d\n", square(5));
    printf("%d\n", add(3, 4));
}
\end{code}

編譯時用 \texttt{gcc main.c math.c -o prog} 把兩個源檔一起編譯。

\subsection{檢查函式回傳值}

當函式回傳一個值時，呼叫者常需要檢查這個值是否代表成功或失敗。 很多
標準庫函式約定用 \texttt{0} 代表成功、非 \texttt{0} 代表各種失敗情況。

\begin{example}{c}{(check-return.c) 檢查 scanf 的回傳值}
#include <stdio.h>

int main(void) {
    int x, y;

    // scanf 回傳成功讀入的項目數
    // 成功讀入兩個整數時回傳 2
    if (scanf("%d %d", &x, &y) != 2) {
        printf("輸入錯誤\n");
        return 1;  // 表示程式失敗
    }

    printf("讀到 %d 和 %d\n", x, y);
    return 0;  // 表示程式成功
}
\end{example}

\iosample{
3 5 \\
}{
讀到 3 和 5 \\
}

\texttt{scanf} 回傳它成功讀入的值的個數。 上例中，若輸入 \texttt{3 5}
就回傳 \texttt{2}（兩個都讀到）；若只輸入 \texttt{3}（或輸入非整數的
字串），就回傳較小的數字或 \texttt{0}。 檢查回傳值是寫穩健程式的習慣。

\begin{drills}[label=drills:func-declare]
\item 函式原型跟函式定義差在哪裡？
\item 標頭檔（\texttt{.h}）裡通常放什麼？源檔（\texttt{.c}）裡放什麼？
\item \texttt{scanf("\%d \%d", \&x, \&y)} 若只成功讀到一個整數，回傳
值會是多少？
\end{drills}

\section{不定個數參數：printf 與 scanf}
\label{sec:func-variadic}

上一節的 \texttt{stdio.h} 只宣告了 \texttt{printf}、\texttt{scanf}
等函式的原型；這一節來看其中一個特殊之處：它們的參數個數在編譯時不固定
——程式設計者寫多少個 \texttt{\%d} 就對應多少個值。 這稱為\textbf{變
參函式}（variadic function），C 標準庫通過格式字串（format string）來支
持這個特性。

\subsection{\texttt{printf} 用法}

\begin{code}{c}{printf 的簽名（簡化版）}
int printf(const char *format, ...);
\end{code}

\texttt{const char *format} 是格式字串，\texttt{...} 代表任意多個參數。
\texttt{printf} 解析格式字串，遇到 \texttt{\%} 轉義序列就從剩餘參數裡
取一個值。

\begin{example}{c}{(printf-demo.c) printf 的不同用法}
#include <stdio.h>

int main(void) {
    // 無參數，只印字面量
    printf("Hello\n");

    // 一個參數
    printf("整數：%d\n", 42);

    // 多個參數
    printf("名字：%s，年齡：%d，分數：%.2f\n", "Alice", 25, 95.7);

    // 常見格式指定
    printf("十進制：%d，十六進制：%x，八進制：%o\n", 255, 255, 255);

    // 字元和字串
    printf("字元：%c，字串：%s\n", 'A', "world");

    return 0;
}
\end{example}

\iosample{
}{
Hello \\
整數：42 \\
名字：Alice，年齡：25，分數：95.70 \\
十進制：255，十六進制：ff，八進制：377 \\
字元：A，字串：world \\
}

常見的格式指定符：
\begin{itemize}
\item \texttt{\%d}：十進制整數
\item \texttt{\%x}、\texttt{\%X}：十六進制（小寫/大寫）
\item \texttt{\%o}：八進制
\item \texttt{\%f}、\texttt{\%.2f}：浮點數（\.2f 表示兩位小數）
\item \texttt{\%s}：字串
\item \texttt{\%c}：單個字元
\item \texttt{\%\%}：印出一個 \texttt{\%}
\end{itemize}

\subsection{\texttt{scanf} 用法}

\begin{code}{c}{scanf 的簽名（簡化版）}
int scanf(const char *format, ...);
\end{code}

\texttt{scanf} 的工作類似 \texttt{printf}，但方向相反：它按格式字串
解析\textbf{輸入}，把讀到的值填進後面參數指向的位置（所以參數要用
\texttt{\&} 取位址）。

\begin{example}{c}{(scanf-demo.c) scanf 的基本用法}
#include <stdio.h>

int main(void) {
    int x;
    float f;
    char c;
    char str[100];

    // 讀一個整數
    printf("輸入整數：");
    scanf("%d", &x);
    printf("讀到 %d\n", x);

    // 讀一個浮點數
    printf("輸入浮點數：");
    scanf("%f", &f);
    printf("讀到 %f\n", f);

    // 讀一個字元（跳過空白）
    printf("輸入字元：");
    scanf(" %c", &c);  // 注意前面的空白
    printf("讀到 %c\n", c);

    // 讀一個字串（不含空白）
    printf("輸入字串：");
    scanf("%s", str);
    printf("讀到 %s\n", str);

    return 0;
}
\end{example}

\iosample{
輸入整數：42 \\
}{
讀到 42 \\
}

\texttt{scanf} 的陷阱：
\begin{itemize}
\item \texttt{\%s} 不檢查長度，容易緩衝區溢出。 用 \texttt{\%99s} 限制長
度比較安全。
\item 在格式字串前加空白（如 \texttt{" \%c"}）能跳過前面的換行符；否則讀
字元時容易誤讀輸入的尾端 \texttt{\textbackslash n}。
\item \texttt{scanf} 的回傳值是成功讀入的項目數，可用來檢查輸入是否正確。
\end{itemize}

\keypoint{\texttt{printf} 和 \texttt{scanf} 是變參函式，通過格式字串決定
參數個數；掌握常見格式指定符對正確處理輸入輸出至關重要。}

\begin{drills}[label=drills:func-variadic]
\item \texttt{printf} 怎麼知道 \texttt{...} 裡實際傳了幾個參數？
\item 讀一個字元前，格式字串為什麼常寫成 \texttt{" \%c"} 而不是
\texttt{"\%c"}？
\item \texttt{scanf("\%s", str)} 有什麼安全疑慮？該怎麼改比較安全？
\end{drills}

\section{高級應用與設計模式}
\label{sec:func-advanced}

\subsection{函式指標回調}

函式本身也有位址，所以我們可以宣告指向函式的指標，在執行期才決定要呼
叫哪一個。 這種寫法常用來做回調（callback），或動態切換演算法。 宣告
語法、跟 \texttt{int *f(int)} 的差異、以及用函式指標陣列切換策略，都
會在指標那一章完整展開；這裡先知道這件事存在即可。

\subsection{遞迴最佳實踐}

遞迴優雅但需謹慎。 底線案例和終止條件很關鍵。

\begin{code}{c}{安全的遞迴寫法}
// 好：有明確的終止情況，且每次呼叫都在逼近它
int factorial(int n) {
    if (n <= 1) return 1;  // 終止情況
    return n * factorial(n - 1);
}

// 差：可能無窮遞迴
// int bad_recursion(int n) {
//     return bad_recursion(n + 1);  // 永遠不會停
// }

// 更好：用記憶化避免重複計算
int fib_memo(int n, int memo[]) {
    if (n <= 1) return n;
    if (memo[n] != -1) return memo[n];
    memo[n] = fib_memo(n-1, memo) + fib_memo(n-2, memo);
    return memo[n];
}
\end{code}

\keypoint{遞迴需明確終止。 複雜遞迴用記憶化最佳化。}

\begin{drills}[label=drills:func-advanced]
\item 函式指標跟一般變數指標，指向的東西有什麼不一樣？
\item \texttt{bad\_recursion} 為什麼永遠不會停？
\item \texttt{fib\_memo} 比沒有記憶化的費氏遞迴版本快在哪裡？
\end{drills}

\addcontentsline{toc}{section}{練習解答}
\section*{練習解答}

各節練習的答案收在這裡。 短題只給能一句話核對的結果。

\paragraph{第~\ref{sec:func-return-value}~節}
\begin{enumerate}
\item 複製一份給參數；\texttt{square} 裡的 \texttt{n} 是呼叫端那個值
的複本，跟呼叫端的變數無關。
\item 先命中 \texttt{year \% 400 == 0}（2000 能被 400 整除），回傳
\texttt{1}。
\item 不會算錯。 2000 能被 4 整除也能被 400 整除，不管先判斷哪一個，
最後都會落在「是閏年」的分支；但若判斷順序寫反、又忘記排除「被 100
整除但不被 400 整除」的情況，其他年份（如 1900）就會算錯。
\end{enumerate}

\paragraph{第~\ref{sec:func-void}~節}
\begin{enumerate}
\item 靠函式內部呼叫 \texttt{printf} 的副作用；它沒有 \texttt{return}
任何值。
\item 會。 陣列參數傳的是位址，\texttt{fillZeros} 對 \texttt{arr[i]}
的寫入會直接反映到呼叫者的原陣列。
\item 合理，例如 \texttt{scanf} 本身：它把值寫進參數指到的位址（副作
用），也回傳成功讀入的個數（回傳值）。
\end{enumerate}

\paragraph{第~\ref{sec:func-params}~節}
\begin{enumerate}
\item 因為閏年的 2 月有 29 天，會影響 3 月以後每個月累加的天數基準。
\item 7。
\item 因為勾股定理要用「最長邊的平方」對「另外兩邊平方和」比較，先找
出哪一邊最長，才知道該把哪個變數放在等號左邊。
\end{enumerate}

\paragraph{第~\ref{sec:func-benefits}~節}
\begin{enumerate}
\item 不是。 它會修改全域變數 \texttt{count}，同樣的參數 \texttt{n}
在不同次呼叫會讓程式的其他狀態（\texttt{count}）改變。
\item 因為 \texttt{min} 不會偷改 \texttt{a}、\texttt{b}，所以
\texttt{minOfThree} 可以放心呼叫兩次 \texttt{min} 而不用擔心中途結果
被污染。
\item 抽象／重用原則：\texttt{lcm} 只負責呼叫 \texttt{gcd} 並組合結
果，相信 \texttt{gcd} 已經正確實作，不需要重新展開它的邏輯。
\end{enumerate}

\paragraph{第~\ref{sec:func-callstack}~節}
\begin{enumerate}
\item 因為整數參數是傳值呼叫：\texttt{swap} 交換的是它自己 stack
frame 裡的複本 \texttt{a}、\texttt{b}，跟 \texttt{main} 的 \texttt{x}、
\texttt{y} 是不同的記憶體。
\item 函式 \texttt{return} 的時候，它的 stack frame 被整個銷毀。
\item 不是。 兩個 \texttt{x} 分別在各自函式的 stack frame 裡，是兩塊
互不相干的記憶體，只是剛好同名。
\end{enumerate}

\paragraph{第~\ref{sec:func-passing}~節}
\begin{enumerate}
\item 整數參數傳的是值的複本；陣列參數傳的是陣列起始位址（第一個元素
的指標）。
\item 會改變；因為函式拿到的是原陣列的位址，對 \texttt{arr[i]} 的寫入
就是在寫原陣列。
\item 因為陣列參數退化成指標後，函式內部無法從 \texttt{arr} 本身得知
陣列有幾個元素，必須額外告知。
\end{enumerate}

\paragraph{第~\ref{sec:func-pitfalls}~節}
\begin{enumerate}
\item 因為 \texttt{n} 是整數參數，只拿到呼叫端那個值的複本；改
\texttt{n} 只改到這份複本，不會回寫到呼叫端的變數。
\item 量到的是指標的大小（通常是 4 或 8 個位元組），不是陣列實際的元
素個數，因為陣列參數在函式簽名裡已經退化成指標。
\item 該。 例如 \texttt{fillWithZeros}／\texttt{reverse} 是好名字，
清楚暗示會修改參數；\texttt{getMax} 這種名字如果偷偷改了陣列，就會誤
導呼叫者。
\end{enumerate}

\paragraph{第~\ref{sec:func-declare}~節}
\begin{enumerate}
\item 函式原型只宣告簽名（名字、參數型別、回傳型別），沒有本體；函式
定義則包含完整的程式邏輯本體。
\item 標頭檔通常放函式原型和型別定義；源檔放實際的函式實作。
\item 1（只成功讀到 \texttt{x} 那一個整數）。
\end{enumerate}

\paragraph{第~\ref{sec:func-variadic}~節}
\begin{enumerate}
\item 靠格式字串：\texttt{printf} 解析格式字串裡每個 \texttt{\%} 轉義
序列，依序從 \texttt{...} 裡取一個對應的參數，並不會另外檢查參數個數
是否吻合。
\item 因為前面的輸入（例如上一個 \texttt{scanf("\%d", ...)}）常留下一
個換行符在輸入緩衝區裡；格式字串開頭的空白會先跳過所有空白字元，避免
直接讀到那個殘留的換行。
\item \texttt{\%s} 不檢查目的地緩衝區的長度，輸入太長會造成緩衝區溢
出；應改用 \texttt{\%99s} 這種帶寬度限制的寫法。
\end{enumerate}

\paragraph{第~\ref{sec:func-advanced}~節}
\begin{enumerate}
\item 函式指標指向一段可執行的程式碼（某個函式），一般變數指標指向資
料（例如一個整數或陣列元素）。
\item 因為它沒有終止情況：每次呼叫都把 \texttt{n} 加大再繼續呼叫自
己，沒有任何一步會停下來回傳。
\item 快在避免重複計算：一般遞迴費氏數列會重複展開同樣的子問題很多
次，記憶化把算過的結果存進 \texttt{memo}，之後遇到同樣的 \texttt{n}
直接查表回傳。
\end{enumerate}

